\begin{figure}[p]
\centering
% --- Knowledge base construction pipeline (TikZ) ---
\begin{tikzpicture}[
  font=\small,
  node distance=8mm,
  box/.style={draw, rounded corners, align=center, inner sep=6pt, text width=0.85\textwidth},
  arrow/.style={-{Latex[length=2mm]}, thick}
]
\node[box] (source) {
\textbf{Source materials} \\
Ariadne Labs Essential Communications Toolkit documents
};
\node[box, below=of source] (parse) {
\textbf{Chunking and structuring} \\
Separate Navigation Map, Conversation Guides, Stuck Points Framework
};
\node[box, below=of parse] (embed) {
\textbf{Embedding generation} \\
Vector representations for semantic retrieval
};
\node[box, below=of embed] (store) {
\textbf{Vector storage} \\
Indexed knowledge base for context retrieval
};
\node[box, below=of store] (retrieve) {
\textbf{Retrieval at inference} \\
Query matching against knowledge chunks
};
\node[box, below=of retrieve] (context) {
\textbf{Context assembly} \\
Retrieved chunks + system prompt $\rightarrow$ LLM
};
\draw[arrow] (source) -- (parse);
\draw[arrow] (parse) -- (embed);
\draw[arrow] (embed) -- (store);
\draw[arrow] (store) -- (retrieve);
\draw[arrow] (retrieve) -- (context);
\end{tikzpicture}
\caption{
Knowledge base construction and retrieval pipeline.
The system used retrieval-augmented generation to ground responses in validated toolkit resources.}
\label{fig:preprocessing_pipeline}
\end{figure}
